\section{Refinement Procedure and Budget Ablation}
\label{app:refine-ablation}

Algorithm~\ref{alg:inv-repair} gives the verifier-guided refinement procedure of Section~\ref{sec:refinement}, and Table~\ref{tab:refine-strategy} gives the category-to-strategy mapping used by its strategy-driven layer. Fig.~\ref{fig:refine-ablation} then isolates the effect of refinement budget on gpt-5.4-mini. The cumulative \emph{Valid} rate rises from \(24.2\%\) at \(k{=}0\) to \(60.2\%\) at \(k{=}3\), \(73.0\%\) at \(k{=}5\), and \(79.2\%\) at \(k{=}10\), and then reaches \(93.3\%\) after the final Houdini pass. The marginal gain after \(k{=}5\) is smaller, while Houdini contributes a final \(+14.1\)\,pp by pruning failing clauses from the best refinement snapshots. In contrast, the \emph{trivial refine} ablation, which uses only a plain error-plus-code prompt, reaches \(70.0\%\) under the same 10-iteration budget and final Houdini pass. This shows that differential refinement prompts and best rollback together account for roughly 23\,pp of additional validity.

\begin{algorithm}[t]
\caption{Verifier-Guided Refinement $\mathrm{Refine}$}
\label{alg:inv-repair}
\small
\KwIn{$\mathcal{F}$ with candidate annotation  $\mathcal{A}_{0}$, symbolic context $\Gamma$, target kind $k \in \{\textit{loop},\textit{spec}\}$, \textit{LLM}}
\KwOut{$\mathcal{F}^*$ with verifier-checked annotation $\mathcal{A}^*$}

$\mathcal{A} \gets \mathcal{A}_{0}$\;
$\mathcal{F}^{-} \gets \mathcal{F} - \mathcal{A}_{0}$\;

\While{$\mathit{Errors} \gets \neg \mathrm{Verify}(\mathcal{F})$}{
    \If{reach iteration limit}{
        \tcp{\small \textit{Layer 3: Houdini-style pruning}}
        $\mathcal{A} \gets \mathrm{Houdini}(\mathcal{F}^{-},\mathcal{A},\mathit{Errors})$\;
        $\mathcal{F} \gets (\mathcal{F}^{-},\mathcal{A})$\;
        $\textbf{break}$\;
    }
    \uIf{$\mathit{Errors} = \textit{syntax error}$}{
        \tcp{\small \textit{Layer 1: syntax repair}}
        $\mathcal{A} \gets \mathrm{RepairCallLLM}(\mathcal{A},\Gamma,\mathit{Errors})$\;
    }
    \Else{
        \tcp{\small \textit{Layer 2: strategy-driven refinement or regeneration}}
        $\mathit{guidance} \gets \mathrm{BuildGuide}(\mathrm{Classify}(\mathit{Errors},k),\Gamma)$\;
        $\mathcal{A} \gets \mathrm{RefineCallLLM}(\mathcal{A},\Gamma,\mathit{guidance})$\;
    }
    $\mathcal{F} \gets (\mathcal{F}^{-},\mathcal{A})$\;
}

\Return $\mathcal{F}$\;
\end{algorithm}

\begin{table*}[t]
\centering
\setlength{\tabcolsep}{5pt}
\caption{Refinement strategies by verifier-failure category.}
\label{tab:refine-strategy}
\begin{tabular}{@{}p{0.14\textwidth}p{0.38\textwidth}p{0.40\textwidth}@{}}
\toprule
\textbf{Target} & \textbf{Verifier feedback} & \textbf{Refinement strategy} \\
\midrule
\multirow{4}{*}{Loop invariant}
 & Establishment (base) obligation fails & Drop or guard facts not implied by the pre-loop state. \\
 & Preservation obligation fails & Rebuild the inductive relation from the loop body's one-step effect. \\
 & Invariant holds but the downstream goal or assertion is unproved & Add loop-exit consequences linking the negated guard to the inductive and unchanged facts. \\
 & Several invariants fail or the candidate is globally inconsistent & Discard the candidate and regenerate from the original template. \\
\midrule
\multirow{4}{*}{Function specification}
 & Target postcondition obligation fails & Drop over-strong claims or split them into path-sensitive clauses. \\
 & Callee-instance obligation fails at a call site & Strengthen caller-side facts before the call, or weaken an over-strong callee precondition. \\
 & Postcondition holds but a user assertion still fails & Add the missing bridge facts to the specification or to the loop-exit facts. \\
 & Frame (\texttt{assigns}) obligation fails & Add only the missing externally visible locations. \\
\bottomrule
\end{tabular}
\end{table*}

\begin{figure}[!htbp]
\centering
\begin{tikzpicture}[x=0.50cm,y=0.056cm,font=\footnotesize,every node/.style={text=figgraytext}]
  % Ours: full \toolname\ refine pipeline (SE-driven context, structured prompts,
  % rollback, ACSL fixer); trivial: minimal "fix-it" prompt, no rollback,
  % no ACSL fixer, no SE context. Both run with refine budget 10 + final Houdini.
  \def\za{24.2}  \def\zb{60.2}  \def\zc{73.0}  \def\zd{79.2}  \def\ze{93.3}
  \def\ta{5.8}   \def\tb{31.0}  \def\tc{40.5}  \def\td{47.0}  \def\te{70.0}
  \def\xH{12.5}
  % horizontal gridlines (faint)
  \foreach \y in {0,20,40,60,80,100} {
    \draw[figgrid] (0,\y) -- (\xH+0.5,\y);
  }
  % axes
  \draw[->,figaxis] (0,0) -- (\xH+0.7,0);
  \draw[->,figaxis] (0,0) -- (0,108);
  \node[anchor=west] at (\xH+0.8,4) {$k$};
  \node[anchor=south] at (0,108) {Valid (\%)};
  % y ticks (skip "0" to avoid colliding with the x=0 tick label)
  \foreach \y in {20,40,60,80,100} {
    \draw[figaxis] (-0.1,\y) -- (0,\y) node[left,xshift=-1pt] {\y};
  }
  % x ticks
  \foreach \x/\lab in {0/0,3/3,5/5,10/10,\xH/{+Houdini}} {
    \draw[figaxis] (\x,-2) -- (\x,0) node[below,yshift=-1pt] {\lab};
  }
  % Ours curve
  \draw[figcontractstroke,ultra thick] (0,\za) -- (3,\zb) -- (5,\zc) -- (10,\zd) -- (\xH,\ze);
  \foreach \x/\y in {0/\za,3/\zb,5/\zc,10/\zd,\xH/\ze} {
    \fill[figcontractstroke] (\x,\y) circle (2.6pt);
    \draw[figcontractstroke,line width=0.8pt] (\x,\y) circle (2.6pt);
    \node[anchor=south] at (\x,\y+2.6) {\y};
  }
  % trivial refine curve
  \draw[figloopstroke,ultra thick,dashed] (0,\ta) -- (3,\tb) -- (5,\tc) -- (10,\td) -- (\xH,\te);
  % all trivial labels above each point with consistent offset
  \foreach \x/\y in {0/\ta,3/\tb,5/\tc,10/\td,\xH/\te} {
    \fill[figloopstroke] (\x,\y) circle (2.3pt);
    \draw[figloopstroke,line width=0.8pt] (\x,\y) circle (2.3pt);
    \node[anchor=south] at (\x,\y+3.2) {\y};
  }
  % legend (lower-right, framed)
  \draw[figaxis,fill=white] (7.0,6) rectangle (12.4,24);
  \draw[figcontractstroke,ultra thick] (7.3,19) -- (8.3,19);
  \fill[figcontractstroke] (7.8,19) circle (2.6pt);
  \draw[figcontractstroke,line width=0.8pt] (7.8,19) circle (2.6pt);
  \node[anchor=west] at (8.5,19) {Ours};
  \draw[figloopstroke,ultra thick,dashed] (7.3,11) -- (8.3,11);
  \fill[figloopstroke] (7.8,11) circle (2.3pt);
  \draw[figloopstroke,line width=0.8pt] (7.8,11) circle (2.3pt);
  \node[anchor=west] at (8.5,11) {trivial refine};
\end{tikzpicture}
\caption{Cumulative \emph{Valid} rate on gpt-5.4-mini over 400 cases as the refinement budget \(k\) increases, for the full \toolname\ pipeline (\emph{Ours}) versus a \emph{trivial refine} baseline.}

\label{fig:refine-ablation}
\end{figure}
